% Chapter 6. Revisit each research question when its experiment finishes; keep the
% limitations section honest and specific.
\chapter{Discussion}
\label{ch:discussion}

\section{RQ1: can the LLM design valid GQDs in hexagon coordinates?}

Yes, provided that the code explains its rejections in terms the LLM can act on. Writing hexagon
positions removed the problems that make long SMILES strings hard to write: the LLM produced
connected lists of roughly the requested size, and almost all of its valid designs were new. What
it could not do was see whether a shape is closed-shell. This is a global property of the whole
structure. A shape that looks reasonable can have one more atom on one sublattice than on the
other, and nothing in a single coordinate pair reveals it. In the first run, 66 of 67 failed
attempts were open-shell.

The two aids address this in different ways. The rule in the prompt gives the LLM a way to avoid
the problem, mainly by building point-symmetric shapes, and it raised the first-attempt pass rate.
The fix edits turn the global condition into a short list of local actions, and they made retries
succeed. Together they reached 99\%, which supports H1. The price is less variety, because the
rule pushes the LLM toward a few safe symmetric shapes. For the agent, which rejects repeats, this
price is acceptable, and the evolutionary loop gets its variety from the parents it edits.

The loop showed that validity is not settled by the representation alone. When the LLM edited
existing GQDs rather than designing new ones, other errors appeared: adding a hexagon that was
already present, removing one that was not, proposing a duplicate, or using more edits than
allowed. The repair step in version 2 recovered many of these children. The general lesson is to
divide the work by what each side can check. The LLM works in a representation where a small
change to the shape is a small change to the text, and the code handles global constraints and
returns actionable feedback.

\section{RQ2: can xTB rank GQDs for \NOtwo{} sensing?}

The small-model check (E3) answers part of this question. For physisorption, which is the regime
of every top candidate, xTB adsorption energies are close to DFT. For chemisorption, xTB is
qualitatively right but overbinds by 0.1 to 0.55~eV, and the reversibility window of the score was
placed with this error in mind. For \ce{O2}, xTB is wrong: it binds a molecule that DFT does not
bind at all. The charge transferred to physisorbed \ce{NO2} is too large in xTB, but it follows the
HOMO energy of the GQD (\(r = 0.88\)), which is the right physical trend.

The results of the loop sharpened this question. For all 40 best GQDs, the reversibility and
selectivity factors equal 1, so the score of the best candidates is the xTB charge alone
(\cref{sec:res-best}). The ranking at the top therefore depended on the one quantity that xTB was
suspected to overestimate.

The GPU benchmark (E4) answers the question, and the answer is split. xTB ranks GQDs well by how
strongly they bind \ce{NO2}, \ce{H2O} and \ce{NH3} (\(\rho = 0.74\) to 0.95) and by their
selectivity margin (\(\rho = 0.86\)), so for binding and selectivity H2 is supported. It does not
rank them by charge transfer (\(\rho = -0.11\)), because in DFT physisorbed \ce{NO2} takes almost
no charge at all (0.01~\(e\) against 0.23~\(e\) in xTB). H2 is therefore contradicted for the
response, and with it for the score (\(\rho = -0.16\)). By the rule fixed in advance, the score has
to change. Two further corrections follow from DFT: water, not CO, is the most dangerous
competing gas, and \ce{O2} does not bind, so the air-stability concern raised by the xTB \ce{O2}
energies disappears at this level of theory.

The consequences for the earlier results are direct. The best GQDs of both the calibration and
the evolutionary loops were selected for a large xTB charge on physisorbed \ce{NO2}, and in DFT
this charge is close to zero. They remain GQDs that bind \ce{NO2} selectively and reversibly, but
they are not shown to give a large electronic response. What the loops optimised was, in effect,
an xTB artefact that follows the HOMO energy. The comparison of the mutators (RQ3) is not affected
as a test of search methods, because both mutators optimised the same objective, but its winners
are not validated sensor candidates.

The broken-symmetry check adds a second problem for the same GQDs (\cref{tab:e4-bs}). The GQDs
with the smallest xTB gaps, which include the best GQDs of both pilot runs, have strong open-shell
character in DFT (\(\langle S^2\rangle \approx 1.5\) to 1.7 and a lowering of up to 0.82~eV), even
though each has a Kekul\'e structure. The closed-shell rule of the pipeline is a test of the
graph, and it does not exclude this. Both the closed-shell xTB scores and the RKS-referenced DFT
energies of these GQDs therefore rest on a reference state that is not the lowest one, and the
strong DFT binding of \ce{NO2} on them is mostly an artefact of that reference. The charge result
does not depend on this, since the nearly closed-shell \code{gqd\_88143} also shows no charge on
\ce{NO2}. The check covers four GQDs with one functional at xTB geometries, so it shows that the
problem exists but not where the gap limit lies. How to treat open-shell character is open
(\code{RESEARCH\_OPTIONS.md}, D22).

Three ways to repair the score are open (\code{RESEARCH\_OPTIONS.md}, D7). The response could be
taken from DFT for the final candidates only, which is accurate but too slow for the loop. It could
be replaced by an orbital descriptor, such as the alignment of the GQD's HOMO with the singly
occupied orbital of \ce{NO2}, which can be computed cheaply and has a clear physical meaning. Or the
target could move to the regime that DFT favours: a weak, reversible covalent bond to \ce{NO2}
(about \(-0.3\) to \(-0.6\)~eV), which in DFT still moves about 0.2~\(e\). The two benchmark GQDs
that keep a non-zero DFT score are of this kind. Separately, the sensor model proposed in
\cref{sec:future} would help with all three, because a real chemiresistor responds to the change
in carrier density of the film, and that depends on more than the net charge on one molecule.

\section{RQ3: does the LLM find better GQDs than random mutation?}

The answer depends on what is measured. With three runs per mutator, the LLM's children scored
higher on average in every pair of runs (0.186 against 0.163 over all runs) and chemisorbed
\ce{NO2} less often, so its edits are better than random ones on average. But the best GQD came
from random mutation in two of the three pairs, and the mean best score was 0.351 for random
mutation against 0.344 for the LLM. For the goal of the search, which is the best GQD found within
the budget, the LLM gave no advantage. H3 is supported for the mean score and not supported for the
best score. The pilot, in which the LLM found the best GQD, was within this run-to-run variation.
Several reasons can explain this outcome.

First, the headroom is small. The best starting GQD scored 0.322 and the best GQD found so far
scores 0.356, a gain of about 10\%. Gains of a few hundredths are close to the noise of the score,
which comes from xTB itself and from probing only two adsorption sites. Second, the score
saturates. Once a GQD physisorbs \ce{NO2} with a good margin over CO, the score equals the charge
transfer, and the remaining differences are small changes in a single number. Third, random edits
with automatic repair are a strong local search on this problem. Most single-hexagon edits produce
a valid, similar GQD, and the population keeps the good ones. The MolLEO study reports a similar
observation for drug-like molecules: in some settings the default graph-based mutation performed
about as well as LLM mutation while using fewer LLM queries \cite{wang2024molleo}.
\authornote{Check this statement against the MolLEO paper before submission.} Fourth, the LLM was asked to make
mostly refinements of the best parents, which favours exploitation over exploration. Careful
refinements raise the average, while the best GQD of a run is an outlier, and the wider variety of
random edits produces more outliers. The pilot also showed that the LLM's multi-edit proposals fail
more often than single edits.

The LLM run has value beyond the best score. Every LLM edit comes with a stated reason, which makes
its choices auditable, and the logged reasons and outcomes are the material for fine-tuning a model
that may propose better edits than either mutator. Three runs per mutator are few, so these conclusions should be
checked with more runs. \pending{Comparison with the baseline without evolution (E6).}

\section{RQ4: what do the best GQDs have in common?}

The calibration (E2) showed the trade-off of H4: a smaller gap raises the charge given to \ce{NO2},
but most small-gap GQDs chemisorb it and lose the ability to recover. The best calibration GQDs
balanced the two at gaps of 0.49 to 0.67~eV. Both loops then found GQDs with smaller gaps, 0.26 to
0.43~eV for the ten best, that still physisorb \ce{NO2} at both probed sites, and these scored
higher. All of the best GQDs have 29 to 40 hexagons, bind \ce{NO2} at about \(-0.43\) to
\(-0.52\)~eV and prefer it over the strongest competing gas by at least 0.16~eV. These results support H4 at the xTB level.

DFT changes this picture (E4). The first half of the trade-off, a larger response at a smaller gap,
rests on the physisorbed xTB charge, which DFT does not reproduce. In DFT, the response survives
only where \ce{NO2} forms a bond, and the two benchmark GQDs with a non-zero DFT score form a weak
one (\(-0.31\) and \(-0.41\)~eV). The design target that DFT suggests is therefore a weak,
reversible bond to \ce{NO2} rather than physisorption, so H4 is not supported at the DFT level.

Two further cautions apply. Chemisorption depends on the local site, and only two sites were
probed, so a GQD that physisorbs \ce{NO2} at its centre and its most exposed edge may still
chemisorb it elsewhere. Small-gap GQDs are generally more reactive, but in DFT \ce{O2} did not bind
eleven of the twelve benchmark GQDs, including both pilot bests, so air stability is not a concern
at this level of theory.

\section{Limitations}
\label{sec:limitations}

\paragraph{Model of the sensor.} Every calculation describes one gas molecule on one isolated GQD
in the gas phase at 0~K. The model leaves out the substrate, the electrical contacts, the
coverage of gas molecules, humidity, temperature and the arrangement of GQDs in a real film. The
response is represented by the charge moved to \ce{NO2}, which is a proxy for the change in
conductance and not a prediction of it. Recovery is estimated from the adsorption energy alone.

\paragraph{Sites and configurations.} Only two adsorption sites are probed, and each interfering
gas is tried in one orientation. Optimisations start from one height and can end in a local
minimum. For final candidates, all symmetry-distinct sites should be checked.

\paragraph{Level of theory.} The score is computed with xTB, whose weaknesses for charge transfer
and \ce{O2} are documented above. The DFT checks are single points at xTB geometries with a small
basis set and no counterpoise correction, and atomic charges depend on the partitioning scheme.

\paragraph{Chemical space.} Only pure-carbon, closed-shell GQDs with 20 to 40 hexagons and no holes
are considered, and shapes whose SMILES RDKit cannot re-read are excluded. ``New'' means absent from
our dataset of 91{,}174 structures, not absent from the literature.

\paragraph{The LLM.} The LLM is a commercial model accessed through an API. Its replies cannot be
seeded, and the model behind a deployment may change over time. The logged prompts and replies make
every run auditable, but an LLM run cannot be repeated exactly. LLM calls are not counted in the xTB
budget, and their cost and time were not measured.

\paragraph{Statistics and design choices.} Three runs per mutator give limited statistical power
(\cref{sec:e5}), and children within a run are not independent. Several settings of the loop, such
as the gap window, the family cap, the amount of history shown to the LLM and the strategy hint in
the prompt, were chosen once and not tuned for either mutator. A different choice could favour one
of them.

\paragraph{Predictor.} The gap predictor is trained on the dataset distribution and is used only as
a filter. Its window can remove good children, as the pilot's 0.30~eV floor did.

\section{Practical cost}

One full xTB evaluation takes about a minute on the laptop, and inside the loop ten children are
scored in about 3 to 6 minutes, so a ten-generation run takes about an hour. A run used 27 to 34
LLM calls. The DFT benchmark of 12 GQDs needs 162 single points, about 12 hours on one A100, or
roughly one hour per GQD with all gases. The work was split by hardware: the xTB loop ran on the
laptop CPU, which was six times faster for xTB than the shared server CPU, and DFT ran on the
server GPU. This split kept each kind of calculation on the hardware that suits it.
